\section{Lean Formalization}
\label{app:lean-formal-core}

The results of \Cref{sec:core-theorem} are formalized in Lean~4 (version 4.28.0) in the package \texttt{ConfluenceFlat}, which has no external dependencies. The development covers abstract rewriting only; the string and term-rewriting computations are not formalized.

\subsection{Setting}

The basic structure is a rewriting system on an arbitrary type, given by a single relation \texttt{step}. There is no finiteness assumption, and every definition and theorem below quantifies over such systems. A separate structure for finite systems extends it with a finite list of states containing every state that occurs in a step; it is provided for representing concrete examples and is not used by any theorem. Reachability is the inductively defined reflexive--transitive closure of \texttt{step}. Joinability, normal forms, local confluence, elementary peaks (with distinct branches), elementary flatness, confluence, normalization, and unique reachable normal forms are defined exactly as in \Cref{def:basic-notions,def:elementary-peak-flatness,def:elementary-holonomy-status}. Termination is well-foundedness of the relation ``$b$ is a one-step successor of $a$'', so that well-founded induction proceeds along rewrite steps. Uniqueness of reachable normal forms means ``at most one''; existence and uniqueness together is a separate predicate.

\subsection{Correspondence}

\Cref{tab:lean-correspondence} lists the Lean theorem corresponding to each result. All names are in the namespace \texttt{ConfluenceFlat.FiniteARS}, a name retained for continuity although the theorems are not restricted to finite systems.

\begin{table}[h]
\caption{Results of \Cref{sec:core-theorem} and their Lean counterparts.}
\label{tab:lean-correspondence}
\centering
\small
\begin{tabular}{@{}ll@{}}
\toprule
Result & Lean theorem \\
\midrule
\Cref{lem:joinable-refl} & \texttt{joinable\_refl} \\
\Cref{thm:elementary-flatness-iff-local-confluence} & \texttt{elementaryFlat\_iff\_localConfluent} \\
\Cref{prop:confluence-normalization-unique-normal-form} & \texttt{existsUniqueNormalFormFrom\_of\_confluent\_normalizing} \\
\Cref{lem:terminating-normalizing} & \texttt{normalizing\_of\_terminating} \\
\Cref{thm:newman} & \texttt{confluent\_of\_terminating\_localConfluent} \\
\Cref{cor:terminating-confluence-iff-flat} & \texttt{confluent\_iff\_elementaryFlat\_of\_terminating} \\
\Cref{thm:normalizing-confluence-iff-unf} & \texttt{confluent\_iff\_uniqueNormalForms\_of\_normalizing} \\
\Cref{cor:terminating-flat-unique-nf} & \texttt{existsUniqueNormalFormFrom\_of\_terminating\_elementaryFlat} \\
\bottomrule
\end{tabular}
\end{table}

Supporting lemmas include \texttt{reachable\_eq\_of\_normalForm} (a normal form reaches only itself), \texttt{eq\_of\_joinable\_normalForms}, and \texttt{uniqueNormalFormFrom\_of\_confluent}. The two directions of \Cref{thm:elementary-flatness-iff-local-confluence} are \texttt{elementaryFlat\_of\_localConfluent} and \texttt{localConfluent\_of\_elementaryFlat}.

\subsection{Axioms}

There are no \texttt{sorry} placeholders and no custom axioms. According to Lean's \texttt{\#print axioms}, the formal versions of \Cref{prop:confluence-normalization-unique-normal-form}, \Cref{thm:newman}, and \Cref{thm:normalizing-confluence-iff-unf}, as well as the supporting lemma that confluence gives unique reachable normal forms, depend on no axioms at all. The formal versions of \Cref{thm:elementary-flatness-iff-local-confluence}, \Cref{lem:terminating-normalizing}, \Cref{cor:terminating-confluence-iff-flat}, and \Cref{cor:terminating-flat-unique-nf} depend only on Lean's standard axioms \texttt{propext}, \texttt{Classical.choice}, and \texttt{Quot.sound}, which enter through classical case distinctions, such as whether two states are equal or whether a state has a successor.
